%% ─────────────────────────────────────────────────────────────────────────────
\section{Problem Description}
%% ─────────────────────────────────────────────────────────────────────────────

A machine paraphrase preserves \emph{what} a text says while rewriting \emph{how} it says it.
Applied repeatedly, this poses the Ship of Theseus question for authorship: after a document
has been paraphrased three times by a large language model, whose text is it?
Authorship attribution, machine-generated text (MGT) detection, plagiarism analysis, and
forensic provenance all assume a stable stylistic signal survives in a document; paraphrasing
cheaply defeats that assumption.

We therefore measure identity decay \emph{quantitatively} rather than as a single pass/fail
outcome, tracking generation by generation which linguistic properties are destroyed first and
which persist. Three questions organize the work:
\textbf{(RQ1)} Which stylometric features decay before semantic content does?
\textbf{(RQ2)} At which generation does authorship attribution become unreliable?
\textbf{(RQ3)} Does the paraphraser leave a recoverable fingerprint, so that authorial identity
is not merely erased but \emph{replaced}?
Our contribution is a decay curve that separates human-author signal loss from model-fingerprint gain.
